% настройки polyglossia
\setdefaultlanguage{russian}
\setotherlanguage{english}


% основной шрифт документа
\setmainfont{Times New Roman}

% перечень использованных источников
% \addbibresource{refs.bib}

% настройка ссылок и метаданных документа
\hypersetup{unicode=true,colorlinks=true,linkcolor=red,citecolor=green,filecolor=magenta,urlcolor=cyan,
	pdftitle={\docname},
	pdfauthor={\studentname},
	pdfsubject={paks\_master},
	pdfcreator={\studentname},
	pdfproducer={Overleaf},
	pdfkeywords={paks\_master}
}

% настройка подсветки кода и окружения для листингов
\usemintedstyle{colorful}
\newenvironment{code}{\captionsetup{type=listing}}{}

% шрифт для листингов с лигатурами
\setmonofont{FiraCode-Regular.otf}[
	SizeFeatures={Size=10},
	Path = templates/,
	Contextuals=Alternate
]

% оформления подписи рисунка
\captionsetup[figure]{labelsep = period}

\floatname{listing}{Листинг}

% подпись таблицы
\DeclareCaptionFormat{hfillstart}{\hfill#1#2#3\par}
\captionsetup[table]{format=hfillstart,labelsep=newline,justification=centering,skip=-10pt,textfont=bf}

% путь к каталогу с рисунками
\graphicspath{{fig/}}

% Внесение titlepage в учёт счётчика страниц
\makeatletter
\renewenvironment{titlepage} {
	\thispagestyle{empty}
}
\makeatother

\counterwithin{figure}{section}
\counterwithin{table}{section}

\titlelabel{\thetitle.\quad}

% для удобного конспектирования математики
\mathtoolsset{showonlyrefs=true}
\theoremstyle{plain}
\newtheorem{theorem}{Теорема}[section]
\newtheorem{proposition}[theorem]{Утверждение}
\theoremstyle{definition}
\newtheorem{corollary}{Следствие}[theorem]
\newtheorem{problem}{Задача}[section]
\theoremstyle{remark}
\newtheorem*{nonum}{Решение}

\def\figureautorefname{Рисунок}

% настоящее матожидание
\newcommand{\MExpect}{\mathsf{M}}

% объявили оператор!
\DeclareMathOperator{\sgn}{\mathop{sgn}}

% перенос знаков в формулах (по Львовскому)
 \renewcommand*{\hm}[1]{#1\nobreak\discretionary{} {\hbox{$\mathsurround=0pt #1$}}{}}

\mathtoolsset{showonlyrefs=false}

\numberwithin{equation}{subsection}

\pagestyle{fancy}
\fancyhf{}

% фиксируем ширину под текст
\setlength{\headwidth}{\textwidth}

% номер страницы
\fancyhead[R]{\thepage}

% убираем линию
\renewcommand{\headrulewidth}{0pt}

% расстояние между header и текстом
\setlength{\headsep}{6mm}

\fancypagestyle{plain}{
	\fancyhf{}
	\setlength{\headwidth}{\textwidth}
	\fancyhead[R]{\thepage}
	\renewcommand{\headrulewidth}{0pt}
}
